\clearpage
\section{Design}

\subsection{Overview}

The three implementations are intentionally separate:

\begin{itemize}[leftmargin=*]
  \item the Virtual Threads version uses imperative recursion and a task per
  directory;
  \item the Reactive version models the scan as a stream of files and uses an
  immutable accumulator;
  \item the Event-Loop version uses Vert.x callbacks and asynchronous file
  system operations.
\end{itemize}

The CLI and GUI can select the implementation at runtime. Both user interfaces
use the same report model and the same conceptual callback contract, while the
scan engine changes underneath. The CLI renders reports as text; the GUI
renders the same data as labels and a table.

\subsection{Library API}

The asynchronous operation \texttt{getFSReport} is exposed by the three
implementation classes:

\begin{itemize}[leftmargin=*]
  \item \texttt{EventLoopFSStat.getFSReport(...)} starts a Vert.x scan and
  returns an \texttt{FSReportJob} cancellation handle;
  \item \texttt{ReactiveFSStat.getFSReport(...)} returns an
  \texttt{Observable<FSReport>}, i.e. a stream of partial and final reports;
  \item \texttt{VirtualThreadsFSStat.getFSReport(...)} starts a virtual-thread
  scan and returns an \texttt{FSReportJob} cancellation handle.
\end{itemize}

The imperative versions deliver intermediate updates, completion, and errors
through an \texttt{FSReportListener}. The reactive version delivers the same
logical information through \texttt{onNext}, \texttt{onComplete}, and
\texttt{onError}. The API therefore follows the conventions of each paradigm
while preserving the same logical result: a sequence of progress reports and a
final filesystem statistics report.

\begin{table}[h]
\centering
\renewcommand{\arraystretch}{1.2}
\begin{tabular}{>{\raggedright\arraybackslash}p{2.7cm} >{\raggedright\arraybackslash}p{3.8cm} >{\raggedright\arraybackslash}p{3.8cm} >{\raggedright\arraybackslash}p{3.8cm}}
\toprule
\textbf{Approach} & \textbf{State ownership} & \textbf{Scheduling model} & \textbf{Main trade-off} \\
\midrule
Virtual Threads & Shared counters with lock-free atomics and a task counter & Recursive task submission on a virtual-thread executor & Simple imperative code, but many tasks for very large trees \\
Reactive & Immutable accumulator passed through the stream & Sequential observable generation with sampled updates & Clean functional model, but no parallel directory walk \\
Event-Loop & Shared asynchronous state guarded by atomics and a visited-set & Vert.x callbacks with asynchronous I/O and periodic updates & Non-blocking design, but more coordination overhead \\
\bottomrule
\end{tabular}
\caption{High-level comparison of the three implementations.}
\end{table}

\subsection{Virtual Threads Strategy}

The Virtual Threads implementation uses a dedicated virtual-thread executor.
Traversal is recursive: every subdirectory is scheduled as a separate task.
This keeps the code close to a normal recursive walk while allowing many
blocking filesystem operations to coexist cheaply. Completion is tracked with:

\begin{itemize}[leftmargin=*]
  \item a shared cancellation flag;
  \item an \texttt{AtomicInteger} task counter;
  \item a \texttt{CountDownLatch} that releases the waiting thread when all work
  is done;
  \item a periodic reporting thread that emits intermediate snapshots.
\end{itemize}

The implementation uses canonical paths to avoid revisiting the same directory
more than once. File sizes are measured through \texttt{File.length()} and each
regular file is assigned to one band through \texttt{FSReport.getBandIndex()}.

Counters are shared because many virtual-thread tasks may discover files at the
same time. The total count and band counters use \texttt{LongAdder}, which is a
good fit for frequent concurrent increments. A simple \texttt{long} would lose
updates, while a single heavily contended atomic counter would scale less well
under many concurrent increments.

\subsection{Reactive Strategy}

The Reactive version uses RxJava 3 and follows a functional style.

The directory traversal is produced by an \texttt{Observable<File>}. The stream
is then transformed into a sequence of immutable scan states with
\texttt{scan(initial, ...)}. Each emitted state contains:

\begin{itemize}[leftmargin=*]
  \item the directory being scanned;
  \item the threshold and the number of bands;
  \item the current file count;
  \item the current histogram;
  \item the start time, used to compute elapsed time.
\end{itemize}

The stream is sampled every 100 ms so that user interfaces are not flooded with
too many updates. Cancellation is handled by disposing the subscription.

This implementation deliberately keeps the directory walk sequential inside the
observable generator. The goal is to keep the reactive control flow explicit
and simple, not to emulate the virtual-thread scheduling strategy.

Because RxJava completion does not carry the last value, the CLI and GUI keep
the latest emitted \texttt{FSReport}. When the stream completes, that report is
forwarded as the final result through the same listener abstraction used by the
other implementations.

\subsection{Event-Loop Strategy}

The Event-Loop version uses Vert.x to keep the traversal asynchronous.

The scan starts by creating a fresh \texttt{Vertx} instance and a periodic
timer that produces progress updates. Blocking filesystem checks, such as
canonical path validation, are wrapped in \texttt{executeBlocking} so they do
not block the event loop.

The asynchronous state tracks:

\begin{itemize}[leftmargin=*]
  \item the number of outstanding tasks;
  \item a cancellation flag;
  \item a set of canonical paths already visited;
  \item the current totals and band counters;
  \item the timer identifier used for progress updates.
\end{itemize}

The directory tree is explored with asynchronous \texttt{readDir()} and
\texttt{props()} calls. Completion is detected when the outstanding task
counter reaches zero.

Unlike the virtual-thread version, the event-loop implementation must be
careful not to block the event loop with canonical path checks. For this reason
validation is executed with \texttt{executeBlocking}. The design also guards
Vert.x shutdown with an atomic flag so that cancellation and normal completion
cannot close the same instance twice.

\subsection{Cancellation and Progress Updates}

The three implementations expose the same user-visible behaviour, but
cancellation is implemented differently:

\begin{itemize}[leftmargin=*]
  \item Virtual Threads set a cancellation flag, interrupt the reporter, and
  stop the virtual-thread executor;
  \item RxJava disposes the subscription, after which the generator observes
  \texttt{isDisposed()} and stops emitting files;
  \item Vert.x sets a cancellation flag, cancels the progress timer, and closes
  the Vert.x instance when pending callbacks have drained.
\end{itemize}

Progress is deliberately throttled. Reports are emitted about every 100 ms,
which avoids flooding the console or the Swing Event Dispatch Thread while
still making long scans observable.

\subsection{User Interfaces}

The CLI accepts:

\begin{itemize}[leftmargin=*]
  \item the target directory;
  \item the maximum file size threshold;
  \item the number of bands;
  \item an optional size unit;
  \item an optional paradigm selector (\texttt{vt}, \texttt{rx}, or
  \texttt{loop}).
\end{itemize}

The GUI exposes the same parameters and adds start/cancel controls. Callback
delivery is marshalled back onto the Swing EDT with
\texttt{SwingUtilities.invokeLater}.

For Virtual Threads and Vert.x, the GUI stores the returned
\texttt{FSReportJob}. For RxJava, it stores the \texttt{Disposable}
subscription. The cancel button checks both handles and cancels whichever one
is active. This keeps the GUI independent from the internal details of the
selected paradigm.
